% AI-written, AI cross-reviewed; no human sign-off; not a claim that RH is proved
% RH Notes blueprint source excerpts for Z2.K8
% Verbatim excerpts in the notation of the parent blueprint. The labelled editions and source ranges appear below.
% See the unit YAML for the current proof status. Earlier-round passages retain their original wording.

% Start of atlas edition annotation. The following block is outside the original excerpt.
\paragraph{Status of named inputs in this reading.}
\noindent\nolinkurl{GlobalTheory.Published.PoitouTateDuality} --- hypothesis at this head (\path{Development/AbelianVarietyTheory/PublishedTheorems/GaloisCohomology₂/ExistsPrimitiveRootPow.lean:13--27}); not compared.\par
\noindent\nolinkurl{GlobalTheory.TateLimit.local_tate_duality_h2_limit} --- in-repo constant (\path{Development/AbelianVarietyTheory/Proofs/FoundationProofs/LengthEqLengthInvariantsOfTorsionLevels.lean:295--302}); not compared.\par
External theorem numbers in the original excerpt are not verified here unless a separate edition-specific source record documents the comparison.
% End of atlas edition annotation.

\section*{Z2.K8. Finite duality and the length formula}

% Parent source SHA-256: 1af64922b98ed2310744411b1f18b25cc426c97220646321471836e1a3a409d6
% blueprint/ch_P.tex lines 1047 through 1098; lem:P-r3-k8
\begin{lemma}[The finite duality sequence at the limit (K8)]\label{lem:P-r3-k8}
The length identity in Proposition~\ref{prop:P-k-main} holds at compact
coefficients, and the cyclotomic character of inertia at $\ell$ over
$\mathbf Q$ maps onto $\mathbf Z_\ell^\times$.
\end{lemma}
\begin{proof}
Take $T_j=(\mathbf Z/\ell^j)(E+1)$ and its Cartier dual $B_j=A_E[\ell^j]$,
$j\ge1$. Write $K_j$ for the kernel of localization on $H^1(G_{L,S},B_j)$
and $D$ for character duality to $\mathbf Q_\ell/\mathbf Z_\ell$.
The order of $T_j$ is a unit in the $S$-integers. The modified local
groups at the infinite places vanish since all are complex, also for
$\ell=2$. Global Poitou--Tate and local Tate duality give
\[
 0\longrightarrow K_j^D\longrightarrow H^2(G_{L,S},T_j)
 \longrightarrow\bigoplus_{w\mid\ell} H^0(L_w,B_j)^D
 \longrightarrow H^0(G_{L,S},B_j)^D\longrightarrow0 .
\]
The first term is the dual of $K_j$ by the global pairing between the
degree-one and degree-two localization kernels. The two right hand maps
are the last two maps of the nine-term sequence. This is the finite
coefficient case of \cite[I.2, Corollary 2.3 and I.4, Theorem 4.10]{S-Milne}.

All terms at finite level are finite since $S$ is finite. Inverse limits
of this sequence are exact: the images in each fixed finite term
stabilize (the Mittag--Leffler condition). The finite sums commute with
the limit. Discrete continuous cochains into $A_E$ have finite image on
each compact power of the group. Taking the filtered union in $j$
therefore commutes with $H^1$, with invariants and with the localization
kernel. Dualizing these direct limits gives the displayed inverse limits.
On the compact side the continuous cochain complex is the inverse limit
of the complexes for $T_j$; cochain reductions are surjective by lifting
each value on its clopen fibre. The possible error in degree two is
$\varprojlim^1 H^1(G_{L,S},T_j)=0$ by the same stabilization. This
identifies the middle group with $H^2(G_{L,S},\mathbf Z_\ell(E+1))$.
Length additivity, and preservation of extended length under duality
as in Lemma~\ref{lem:P-r3-k1}, give the asserted identity. The source
proves the kernel and cokernel length comparisons in \nolinkurl{ptShaAt},
\nolinkurl{ptCokerAt} and the local term in \nolinkurl{local_tate_duality_h2_limit}
(\path{LengthEqLengthInvariantsOfTorsionLevels.lean:462,393,295}).
They enter \nolinkurl{vFormula_of_classicalTheorems}
(\path{GaloisCohomologyLength/Formula.lean:15--27}).

Finally $\Phi_{\ell^j}(1+X)$ is monic, reduces to
$X^{\varphi(\ell^j)}$ modulo $\ell$, and has constant term $\ell$.
By Eisenstein, $\mathbf Q_\ell(\mu_{\ell^j})/\mathbf Q_\ell$ is
totally ramified of degree $\varphi(\ell^j)$. Inertia thus maps onto
$(\mathbf Z/\ell^j)^\times$ for each $j$. Its image in the inverse
limit is closed by compactness and dense by these surjections.
See \nolinkurl{cycZ_Ine_surjective}
(\path{InertiaEigenvalues/CyclotomicCharInertiaSurjective.lean:63--96}).
\emph{Status: complete under the finite duality theorems in the proposition.}
\end{proof}

